% AI Cashier - few-shot product recognition for a low-cost self-checkout.
%
% Tables and figures are generated by research/report.py from research/results/.
% Do not edit anything in tables/ or figures/ by hand - re-run the experiments.
%
% Anything marked \todo{} is a number or a claim that must come from a real
% capture session. Every one of them must be gone before this is submitted.

\documentclass[10pt,twocolumn]{article}
\usepackage[margin=0.85in]{geometry}
\usepackage{booktabs,graphicx,amsmath,microtype,url}
\usepackage[hidelinks]{hyperref}
\usepackage{xcolor}

\newcommand{\todo}[1]{\textcolor{red}{\textbf{[TODO: #1]}}}

\title{\textbf{Teaching a Self-Checkout a New Product in Under a Minute:\\
Few-Shot Retrieval and Multimodal Disambiguation on a \pounds 100 Till}}

\author{Punn, Pleum, Athens, Kit, Bible, Pokpong \\
{\small Assumption College Sriracha, Chonburi, Thailand} \\
{\small Engineering Design and Innovation}}

\date{}

\begin{document}
\maketitle

\begin{abstract}
Camera-based self-checkout removes the barcode, but the usual implementation
replaces it with a worse constraint: a closed-set detector knows exactly the
products it was trained on, so a shop that changes its range must collect a
dataset and retrain. It also cannot abstain --- shown something it has never
seen, a softmax still names a product, and the customer is charged for it.

We rebuild a working closed-set self-checkout as a retrieval system. A
class-agnostic proposer finds objects, a small embedding network describes them,
and identity is a nearest-neighbour lookup in a gallery of reference views.
Adding a product becomes adding rows to that gallery: the operator puts it on the
mat, turns it a few times, and it is on sale. Rejection becomes a distance
threshold, so the till can say it does not know. Because appearance alone cannot
separate same-brand variants that differ mainly in one printed word, we combine
it with mass from a load cell and physical size from a fiducial marker on the mat.

On a Raspberry~Pi~5 the pipeline runs at \todo{FPS from research/bench.py}. Enrolling
a product takes \todo{seconds, E8} against \todo{hours, the team's own record} to
retrain the closed-set model it replaces. Over \todo{N} baskets the mean pricing
error is \todo{THB, E7}. We report the case the method does not solve --- visually
near-identical variants --- and how much of it mass and size recover.
\end{abstract}

\section{Introduction}

A self-checkout that recognises products by sight is attractive for a small shop:
no barcode to find, no scanner to aim, nothing to key in when a label is torn.
Building one with an object detector \cite{redmon2016yolo} is now
straightforward, and our own version~1 did exactly that, recognising twelve
snack and drink lines with two YOLOv8n models.

It stopped being useful the moment the shop wanted a thirteenth product.

This is not an implementation detail. A closed-set detector's output space is
fixed when it is trained. Extending it means photographing the new product many
times, labelling every image, retraining, validating and reinstalling the model
--- days of work for one line of crisps, in a setting where the range changes
weekly. Worse, the classifier cannot represent ignorance: presented with a
product not in its training set, it distributes probability over the classes it
does have, and the till confidently charges the wrong price. In a system that
takes money, a confident wrong answer is more damaging than no answer.

We therefore reformulate recognition as retrieval. Our contributions:

\begin{enumerate}
\item \textbf{A self-checkout that is taught rather than trained.} Enrolling a
product is $k$ reference views appended to a gallery, with no gradient step
anywhere in the loop. We measure the wall-clock cost against retraining
(Section~\ref{sec:results}).
\item \textbf{An abstention mechanism the closed-set system structurally lacks.}
Rejection is a threshold on gallery distance, calibrated on held-out products
rather than chosen by hand, and reported as an open-set detection problem.
\item \textbf{Multimodal disambiguation for the case appearance cannot solve.}
Same-brand variants are near-identical to an embedding; they are not
near-identical on a scale. We fuse appearance, mass and physical size as summed
log-likelihoods and give the per-modality ablation.
\item \textbf{Measurements on the hardware it actually runs on}, including a
negative result: dynamic INT8 quantisation of our backbone was both slower and
substantially less faithful on ARM, so we ship FP32.
\end{enumerate}

We claim application and measurement, not novelty of technique. Retrieval-based
product recognition is established \cite{wei2019rpc}; what is not established is
what it costs and what it buys on a till a school can build.

\section{Related work}

\paragraph{Retail product recognition.} Recognising packaged goods from one
studio reference per product was posed by GroZi-120 in 2007
\cite{merler2007grozi}; the survey of Wei et al.\ \cite{wei2020survey} names
automatic checkout as the first application of the field. Checkout-specific
benchmarks are RPC \cite{wei2019rpc} (200 SKUs, 30k tray scenes), the AI City
Challenge retail track \cite{naphade2022aicity,naphade2023aicity}, which trained
on rendered products and tested on real video, and the 2026 GroceryVision
retrieval challenge \cite{praw2026groceryvision}. Fine-grained SKU corpora
\cite{bai2020products10k,peng2020rp2k,klasson2019grocery,chen2022unitail} and
the flavour-variant benchmark MIMEX \cite{tur2024mimex} document the failure
that packaging alone cannot resolve. Our setting is smaller than any of these
but carries the constraint they omit: the product range changes after the
system is deployed, and the deployer has no data team.

\paragraph{Recognition by retrieval.} Comparing a query with a few labelled
references is standard metric learning
\cite{schroff2015facenet,snell2017protonets,khosla2020supcon,deng2019arcface};
for few-shot classification the backbone matters more than the algorithm
\cite{chen2019closerlook}. Retail-specific embeddings address the shift between
reference and store images \cite{tonioni2019domain}, and detector-plus-embedding
stacks are the design recent reviews recommend when SKU churn is high
\cite{review2025oneshotretail,icoiact2025fewshotretail}. NVIDIA's published
self-checkout reference follows the same decomposition---detect, embed, retrieve
against a reference database, verify against the barcode scanner---and
recommends 20--30 reference images per product \cite{nvidia2024retailrecognition};
Mashgin advertises sub-minute enrolment on more than 10,000 deployed units
\cite{mashgin2015patent,couchetard2022mashgin}. General-purpose features
\cite{radford2021clip,zhai2023siglip,oquab2023dinov2,simeoni2025dinov3} make
retrieval work with no in-domain training; a 2026 study of 190 open encoders on
grocery retrieval finds that training-data quality outweighs model size and that
mobile-scale encoders are competitive \cite{maminta2026grocery}, which is the
result our hardware budget depends on. We deploy an ImageNet-pretrained
MobileNetV3 \cite{howard2019mobilenetv3,deng2009imagenet} and report larger
encoders as a ceiling.

\paragraph{Knowing what you do not know.} Open-set recognition was formalised
by Scheirer et al.\ \cite{scheirer2013openset}; surveys
\cite{geng2020ossurvey,yang2021oodsurvey} fix the vocabulary we use. Maximum
softmax probability \cite{hendrycks2017baseline} and the energy score
\cite{liu2020energy} are the baselines; Vaze et al.\ \cite{vaze2022goodclosedset}
show open-set quality tracks closed-set accuracy, which predicts that our backbone
ablation moves both. A retrieval head has a more natural statistic than any of
these: the distance to the nearest enrolled reference, which is defined for
products the network never saw. Open-world detection \cite{joseph2021owod}
extends the problem to incremental classes without forgetting, the direction
we leave to future work.

\paragraph{Fusion and verification.} Vision plus shelf weight reaches 92.6\%
item identification where vision alone manages 60--70\% \cite{falcao2020faim};
Amazon's store-scale system fuses ceiling cameras with shelf load cells. The
loss literature explains why verification matters commercially: stores with
self-checkout lose 0.42 percentage points more than stores without, and one to
five percent of transactions contain a missed scan \cite{hopkins2026sco}, about
half of it accidental \cite{beck2018sco,taylor2016swipers}. Our weight and size
terms are a verification signal in that sense, at the till rather than the shelf.

\paragraph{Proposals and tracking.} Class-agnostic proposals can come from
foundation models \cite{kirillov2023sam,liu2023groundingdino}; on a fixed mat
under fixed light, background subtraction is two orders of magnitude cheaper and
is the ablation we report. Associating detections across frames
\cite{zhang2022bytetrack} keeps one item from becoming several cart lines; our
objects are stationary, so a centroid tracker with per-track voting suffices.

\section{System}

\paragraph{Hardware.} A Raspberry~Pi~5 (8\,GB) with a 14\,\char`\"{} touchscreen, a
USB camera on a fixed overhead mount, a diffused ring light, a matte mat carrying
a printed ArUco marker of known side length, and a 5\,kg single-point load cell
read through an HX711. Controlled illumination is not cosmetic: it removes most
of the appearance variation the embedding would otherwise have to be invariant to.

\paragraph{Software.} One process. The till (PySide6) owns the camera, the
load cell and the cart and writes SQLite directly; a FastAPI server on a
background thread serves the shopkeeper's dashboard --- inventory, takings and
the deployment log --- as a read-mostly view on the same file, reachable from a
phone on the shop network behind a PIN. Products, stock and sales live in
SQLite, and a sale is one transaction: stock down, sale written, payment closed,
or none of it. The cell is read continuously on its own thread and re-zeroed
whenever the empty pan settles, so drift never carries between baskets.

\subsection{Deployment and compliance}

Three properties follow from the design and are worth stating because a
deployable till is judged on them. \emph{Privacy:} no frame is stored; only
product crops are embedded and only the 576-dimensional vector persists, which
keeps the system inside the ``legitimate interest'' basis Thai practitioners
apply to shop cameras under the PDPA and away from biometric data.
\emph{Restricted goods:} Thai law permits alcohol sales only 11:00--24:00 and
only to verified adults, and forbids tobacco sale by machine or display; the
product table carries a restriction flag and such items are never completed
unattended. \emph{Metrology:} the load cell verifies a basket, it never prices
one, so it is a sensor rather than a trade weighing instrument, which is what
allows a THB-300 cell to be used at all. \todo{cite docs/research/01 rows L03--L06,
L12, L14 in the camera-ready if the venue allows non-academic references}

\section{Method}
\label{sec:method}

\paragraph{Proposal.} Objects are found by differencing the frame against a
photograph of the empty mat, per colour channel rather than in luminance --- a
grey difference is blind to any product whose brightness happens to match the mat.
The proposer is class-agnostic by construction, which is what allows a product
the system has never seen to be located at all.

\paragraph{Embedding.} Each crop is resized to $224^2$ and passed through an
ImageNet-pretrained trunk with its classifier removed, giving a
$d$-dimensional descriptor. The trunk is exported to ONNX so the till never
imports a training framework.

\paragraph{Gallery and centring.} A product is $k$ vectors. Identity is the
nearest reference by cosine similarity. Crucially, similarities are computed
after subtracting a centre: an ImageNet trunk is not trained to be
contrastive, and every pair of natural images shares a large common direction.
On our development set, two entirely different products still scored $0.80$
on average before centring, leaving no room for a threshold; centring moved that
to $-0.13$ while views of one product stayed near $0.92$. The centre is the gallery mean
\emph{frozen} once enough products are enrolled: a centre that floated with
every enrolment would move every existing score and silently invalidate the
threshold below. PCA-whitening fitted on held-out products~\cite{jegou2012negative}
is evaluated as the alternative in the backbone table.

\paragraph{Rejection.} A query whose best similarity falls below $\tau$ is
reported as unknown. $\tau$ is chosen as the lowest value accepting a target
fraction of enrolled products, measured on a validation split --- never fixed by
hand. Anchoring on the true-positive rate rather than on accuracy reflects the
asymmetric cost: refusing a stocked product delays a queue, whereas accepting an
unstocked one takes the wrong money.

\paragraph{Fusion.} Appearance, mass and size each contribute a log-likelihood:
\[
s(c) = \frac{\cos(q, c)}{T}
     - \frac{w_m}{2}\!\left(\frac{\hat{m}-m_c}{\sigma_m}\right)^{\!2}
     - \frac{w_z}{2}\!\sum_{j}\!\left(\frac{\hat{z}_j-z_{c,j}}{\sigma_z}\right)^{\!2}
\]
A modality with no measurement or no reference value contributes exactly zero, so
it never tilts the ranking and never penalises a product for incomplete data.
The mass term is used only when exactly one item was placed since the cart last
changed --- the pan's difference is then one item's mass; with several items on
the mat, appearance and size decide and the basket check below verifies the
total. Rejection stays an appearance decision: agreeing mass must not rescue
something that does not look like anything enrolled.

\paragraph{Basket verification.} Separately from identification, total mass on
the pan is compared with the sum of the basket's reference masses. Independent
errors add in quadrature, so the tolerance is
$k_\sigma\sqrt{\sigma_{\text{cell}}^2 + n\sigma_{\text{item}}^2}$; summing them
linearly gives a band wide enough to swallow the swaps the check exists to catch.
When any item lacks a reference mass the check reports that it could not be
performed rather than that it passed.

\section{Experimental protocol}
\label{sec:protocol}

All recognition results are reported on products held out from any
representation training, enrolled only from $k$ views. Evaluating on products
inside the training set would measure nothing. Splits are produced by
\texttt{research/dataset.py} and recorded in every result file.

\todo{state the final product count, how many are in the legacy detector's
classes, and how many near-identical pairs the set contains}

The experiments are: closed-set baseline (E1); accuracy against $k$ (E2);
backbone against latency (E3); temporal voting (E4); open-set rejection (E5);
fusion ablation and item-swap detection (E6); end-to-end basket pricing error
(E7); and the cost of adding one product (E8).

\section{Results}
\label{sec:results}

\input{tables/e1_closed_set}
\input{tables/e2_fewshot}
\input{tables/e3_backbones}
\input{tables/e4_temporal}
\input{tables/e5_openset}
\input{tables/e6_fusion}
\input{tables/e6_swap}
\input{tables/e6_escalation}
\input{tables/e7_basket}
\input{tables/e8_enrolment}
\input{tables/e9_public}
\input{tables/e9_public_openset}

\begin{figure}[t]
\centering
\includegraphics[width=\columnwidth]{figures/e2_fewshot.pdf}
\caption{Accuracy against the number of reference views, on products the
representation never saw.}
\label{fig:fewshot}
\end{figure}

\begin{figure}[t]
\centering
\includegraphics[width=\columnwidth]{figures/e6_swap.pdf}
\caption{Item-swap detection against weight tolerance. The operating point is a
choice about how often the till may accuse an honest customer.}
\label{fig:swap}
\end{figure}

\paragraph{Public benchmark.} Table~\ref{tab:public} and Table~\ref{tab:publicopenset}
repeat E2 and E5 on the Grocery Store dataset~\cite{klasson2019grocery}: phone
photographs of 81 products on Swedish shop shelves, with one manufacturer
(``iconic'') image per class. There is no mat, so the whole photograph is the
crop; the split, the gallery rule and the threshold rule are exactly those of the
rig experiments, run by \texttt{research/prepare\_grocerystore.py} and
\texttt{run.py --source folder}. The \emph{packages} rows use the 30 carton
classes (milk, juice, yoghurt), the kind of object a till sees; the \emph{iconic}
rows put the manufacturer image first among the reference views, so $k=1$
measures enrolment from a pack shot with no capture at all. These rows say what
the frozen representation does on photographs nobody in the team took; they say
nothing about the rig, its light or its mat. The centred cosine of
Section~\ref{sec:method} beats the raw cosine at every $k$ on these photographs as
it did on the mat.

\paragraph{What that costs.} Table~\ref{tab:backbones} prices it. Every row is
measured under the same runtime the till uses, so the milliseconds are
comparable rather than indicative; the harness reports that condition and
refuses the claim when it does not hold. The accuracy is not free, and the
factor is large: the cheapest CLIP-class encoder here costs an order of
magnitude more per crop than the ImageNet trunk it would replace, and the most
accurate costs nearly two. Whether any of that fits an interactive frame budget
on the target hardware is a separate measurement, and until it exists the
comparison bounds the representation rather than the product. Two details
matter for anyone repeating this. The two S-variants are not ordered as their
names suggest --- the smaller is both more accurate and slightly cheaper, so the
larger is dominated on both axes. And one otherwise reasonable candidate, a
self-supervised transformer, could not be exported to the inference runtime at
all: a representation that cannot be frozen is not a candidate for an edge
device however it benchmarks, which is a criterion the retrieval literature
rarely applies.

\paragraph{The encoder is the lever.} Every row of Table~\ref{tab:public} is the
same matching code, the same split and the same enrolment rule over a different
frozen encoder, so the spread between rows is attributable to the representation
alone. It dominates: exchanging the ImageNet-trained trunk for a small
CLIP-class encoder~\cite{maminta2026grocery} moves top-1 further than any choice of
$k$, of scoring rule, or of tracker in this paper. Two things about that are
worth stating rather than tabulating. First, the gain is larger in abstention
than in identification --- the proportion of unstocked products wrongly accepted
at a fixed true-positive rate falls by an order of magnitude
(Table~\ref{tab:publicopenset}) --- and it is abstention, not top-1, that decides
whether a till takes the wrong money. Second, the two properties come apart: the
self-supervised encoder matches the ImageNet trunk's top-1 on the carton subset
while separating known from unknown considerably better. A representation can be
no better at naming a product and still much better at admitting it has never
seen one, which is a distinction the self-checkout literature reports only
half of.

\todo{write the rest of this section only after research/run.py --source captures. Every
table above regenerates from the results; do not type numbers into the prose
that are not in a table.}

\paragraph{When an unattended till should fetch a person.} A till with nobody
standing at it has to answer a discrepancy on its own, and the useful question is
not whether it can detect a swap but what it should \emph{do} about one.
Table~\ref{tab:escalation} sweeps the value at risk at the shipped weight
tolerance. Detection is constant down the column --- it depends on the tolerance,
not on the price --- so what the sweep isolates is the cost of the response:
raising the line hands more discrepancies back to the customer to resolve and
fewer to staff, and the last column is the bill for the one before it, in honest
baskets interrupted per thousand. Setting that line is therefore an economic
decision rather than a matter of taste, and it is one no commercial system
publishes, because publishing it means publishing a false-alarm rate. We report
ours. The till never acts on a customer physically: the responses are a prompt,
holding the sale, and calling a member of staff, and the reasoning for stopping
there is in Section~\ref{sec:method}'s abstention argument --- most self-checkout
loss is accidental~\cite{beck2018sco,taylor2016swipers}, so most of what a stricter line would catch is
an honest mistake.

\paragraph{Enrolment without capture.} The \emph{iconic} rows enrol each product
from the manufacturer's own product image before any shelf photograph, so $k=1$
on those rows is enrolment with no capture at all. With the ImageNet trunk that
setting is the worst in the table; with the CLIP-class encoder it is among the
best, beating the same encoder enrolled from a single real photograph and
matching it at five. A canonical, evenly-lit product image is a better single
reference than one arbitrary photograph of a shelf, and a representation trained
on web image--text pairs is the thing that can exploit it. If this transfers to
the rig it removes the capture step from enrolment entirely for any product whose
manufacturer publishes an image --- which changes the operator cost the paper
measures in E8, and is the reason that experiment reports operator seconds
separately from compute seconds.

\section{Limitations}

\paragraph{Benchmark transfer.} The public rows are shelf photographs, not mat
crops under a fixed light: they bound what the encoder can do, not what the till
does, and the near-identical cartons in that set (six Arla milks) are the same
failure as ours.

\paragraph{The best encoder is not yet the shipped one.} The strongest row of
Table~\ref{tab:public} comes from a vision transformer an order of magnitude
heavier than the network the till runs, measured on a laptop. Whether it holds an
interactive frame budget on the target hardware is unmeasured, so the result
bounds the representation rather than the deployed system; the distilled variants
exist for exactly this trade and are the obvious next measurement.

\paragraph{Near-identical variants.} Two flavours of the same brand in the same
packaging are the case appearance does not solve, and no amount of enrolment
fixes it --- the information is not in the picture at the resolution we work at.
Mass and size recover part of it; the residual is reported rather than averaged
away. \todo{give the per-pair accuracy}

\paragraph{The scale has a matching blind spot.} Swapping one product for another
of nearly the same mass is invisible to the basket check. The two modalities
cover each other's failures, which is the argument for fusing them, but neither
covers its own.

\paragraph{Controlled conditions.} Everything assumes the fixed rig: one camera
position, one light, a known mat. This is realistic for a counter-top till and
not realistic for a handheld or trolley-mounted system.

\paragraph{Scale of evaluation.} \todo{state the number of products honestly and
say what it does not support --- a few dozen SKUs is not evidence about a
supermarket's long tail}

\paragraph{Simulated adversaries.} Our swap attempts are constructed by us, not
by people trying to defeat the till for gain.

\paragraph{No public-benchmark result yet.} All numbers come from our
own rig. A held-out-SKU split of RPC \cite{wei2019rpc} and the GroceryVision
retrieval task \cite{praw2026groceryvision,maminta2026grocery} are the obvious
external checks and are \todo{done / not done}.

\paragraph{Legal gating is designed, not evaluated.} The age and time gate for
restricted goods and the privacy properties are asserted from the design; we do
not measure how often staff override them or how a shop actually uses them.

\paragraph{The idea is not new; the measurement is.} Learning products at the
till was patented before 2015 and is sold at scale today. Our contribution is
the quantified trade-off on low-cost hardware, and we make no claim beyond it.

\section{Conclusion}

Replacing a closed-set detector with a retrieval head changes what a
camera-based till can do in a shop rather than in a demonstration: the range can
change without an engineer, and the machine can decline to guess. The costs are a
second sensing modality for same-brand variants, and a rig that must stay still.

\bibliographystyle{plain}
\bibliography{refs}

\end{document}
